\documentclass{article}

% Official template: https://neurips.cc/Conferences/2026/CallForPapers
% Default: anonymous Main Track submission with line numbers.
% For a non-anonymous draft, use \usepackage[preprint]{neurips_2026}.
\PassOptionsToPackage{numbers,sort&compress}{natbib}
\usepackage{neurips_2026}

\usepackage[utf8]{inputenc}
\usepackage[T1]{fontenc}
\usepackage{hyperref}
\usepackage{url}
\usepackage{booktabs}
\usepackage{amsmath}
\usepackage{amsfonts}
\usepackage{microtype}

\title{Drift-Regularized Multi-Token Prediction via On-Policy Self-Distillation}

% Author details and affiliation text copied from the supplied manuscript.
% They are hidden in the default anonymous submission mode.
\author{%
  Po-Wei Huang, Feng-Ting Liao \\
  Department of Computer Science, University of Oxford \\
  MediaTek Research
}

% Adapted from the supplied July 2026 manuscript:
% Multi_Token_Prediction_via_Self_Distillation_Drifting_Model_Revision_.pdf
% Its numerical results are explicitly synthetic placeholders; this proposal
% retains the motivation, method, and experimental plan without those results.

\begin{document}

\maketitle

% \begin{abstract}
% Standard autoregressive language models incur substantial latency by generating
%  tokens strictly sequentially. Multi-token prediction (MTP) addresses this by
% predicting a block of future tokens in a single forward pass. However, training
% MTP models solely via token-level cross-entropy leaves a representation gap:
% parallel hidden states lack explicit guidance on how to organize intermediate
% predictive representations. We propose a drift-regularized self-distillation
% framework that constructs a local conditional attraction--repulsion field to
% align student hidden-state populations with teacher feature distributions.
% Unlike paired mean squared error (MSE), this objective does not require
% point-to-point Euclidean matching. We will test whether conditional drift
% improves task accuracy, calibration, latent-space alignment, and the
% accuracy--acceleration trade-off relative to token-only MTP and paired MSE.
% \end{abstract}

\section{objective}
\label{sec:objective}
Standard next-token prediction generates one token per step, incurring high
memory-bandwidth overhead and inference latency during long-form generation and
reasoning-intensive tasks. MTP mitigates this bottleneck by generating a
contiguous block of $k$ tokens in a single forward pass. We build on on-policy
self-distillation, in which a student proposes future tokens at inserted mask
positions and a frozen teacher supplies supervision conditioned on the
student-generated continuation~\citep{kirchenbauer2026mtp}.

Our objective is to reduce the capability loss that accompanies longer parallel
prediction horizons. We will test the hypothesis that, at comparable realized
acceleration, token-only MTP is outperformed by MTP with paired latent MSE,
which is in turn outperformed by MTP with conditional drift. Success will be
assessed through task scores relative to the original sequential teacher,
confidence calibration, held-out feature-distribution diagnostics, and measured
throughput and latency. This ordering is a hypothesis to be tested rather than
an established result.

\section{novelty}
\label{sec:novelty}
A correct token prediction constrains relative logit margins without uniquely
determining the hidden representation. Consequently, newly inserted MTP mask
states must learn parallel future-token computations without explicit
supervision over their representation geometry. Work on feature-based
speculative decoding, including EAGLE~\citep{li2024eagle}, motivates addressing
feature uncertainty; NextLat~\citep{teoh2025nextlatent} also provides motivation
for auxiliary prediction in latent space.

Paired MSE supplies dense supervision but imposes a one-to-one Euclidean match
between each parallel student state and its sequential teacher counterpart.
Our proposed contribution is to adapt population-level drifting
\citep{deng2026drifting} to conditional MTP representation alignment. A local
vector field attracts student states toward teacher feature neighborhoods and
repels them from crowded student neighborhoods. Conditioning interactions on
prediction horizon and prefix context is intended to preserve relevant local
structure. We will compare horizon-wise alignment with block-level alignment
to determine whether modeling inter-token feature dependencies improves chunk
coherence.

\section{impact}
\label{sec:impact}
If successful, this approach could extend the usable parallel prediction horizon
while retaining more of the sequential model's reasoning capability. The
regularizer is applied during training and adds no inference module to the
standalone MTP model. The intended practical benefit is improved generation
quality at a given latency or greater throughput at a given quality target.
These benefits require system measurements and cannot be inferred from the
configured block size alone.

The study will also provide a controlled comparison of output-only supervision,
pointwise feature regression, and population-level feature alignment. Reporting
calibration separately from latent distribution distances will help distinguish
reliable confidence from representation matching. A later extension could apply
the same conditional regularizer to a speculative drafter such as
DFlash~\citep{chen2026dflash}; the initial study will focus on standalone MTP.

\section{Methods}
\label{sec:methods}

\subsection{On-policy MTP self-distillation}
For training region $r$, let $c_r$ denote a context prefix and $k_r$ the number
of future predictions. Following~\citet{kirchenbauer2026mtp}, the flattened
student input inserts $k_r-1$ trainable \texttt{<MTP>} mask tokens. Repeated
position IDs and a blocked causal attention mask allow active prediction states
to attend to the prefix and preceding masks without altering copied context
representations. At horizon $j$, the student produces hidden state $u^S_{r,j}$
and token distribution $p^S_{r,j}$. Its greedy proposal is
$\hat y_{r,j}=\arg\max_v p^S_{r,j}(v)$.

The frozen teacher is evaluated on the student-forced continuation:
\begin{equation}
  u^T_{r,j}=H_{\theta_T}(c_r,\hat y_{r,<j}),
  \qquad
  \bar y_{r,j}=\arg\max_v p_{\theta_T}(v\mid c_r,\hat y_{r,<j}).
\end{equation}
With $N_{\mathrm{act}}=\sum_r k_r$, the hard on-policy distillation loss is
\begin{equation}
  \mathcal L_{\mathrm{MTP}}
  =-\frac{1}{N_{\mathrm{act}}}\sum_{r,j}
    \log p^S_{r,j}(\bar y_{r,j}).
\end{equation}
For active index $i=(r,j)$, let $x_i=\phi_S(u^S_i)$ and
$y_i^+=\phi_T(u^T_i)$ be features in a shared latent space, using the same
normalization. Detached student features $y_l^-=\operatorname{sg}(x_l)$ form
the negative population, where $\operatorname{sg}$ denotes stop-gradient.

\subsection{Latent regularization}
We compare the unified objective
\begin{equation}
  \mathcal L=\mathcal L_{\mathrm{MTP}}+\lambda_R\mathcal L_R,
  \qquad R\in\{\mathrm{MSE},\mathrm{drift}\}.
\end{equation}
The paired regression baseline is
\begin{equation}
  \mathcal L_{\mathrm{MSE}}
  =\frac{1}{N_{\mathrm{act}}}\sum_i
    \left\|x_i-\operatorname{sg}(y_i^+)\right\|_2^2.
\end{equation}
For conditional drift, teacher and detached student pools $\mathcal I_i^+$ and
$\mathcal I_i^-$ interact through
\begin{equation}
  \kappa_{\mathrm{cond}}(i,l;y_l)
  =\kappa_h(x_i,y_l)\,
    \kappa_c(g(c_i),g(c_l))\,\mathbf 1[a_i=a_l],
\end{equation}
where $a_i=j$ identifies the prediction horizon, $g(c_i)$ is a detached prefix
embedding, and $\kappa_h$ and $\kappa_c$ are nonnegative distance kernels over
features and context, respectively. For $s\in\{+,-\}$, define
$\mathcal J_i^+=\mathcal I_i^+$ and
$\mathcal J_i^-=\mathcal I_i^-\setminus\{i\}$. Normalize affinities as
\begin{equation}
  a^s_{il}
  =\frac{\kappa_{\mathrm{cond}}(i,l;y_l^s)}
  {\sum_{m\in\mathcal J_i^s}\kappa_{\mathrm{cond}}(i,m;y_m^s)+\epsilon},
  \qquad l\in\mathcal J_i^s,
\end{equation}
with $\epsilon>0$ for numerical stability. The self-excluded drift estimate is
\begin{equation}
  \widehat V_i
  =\sum_{l\in\mathcal J_i^+}a^+_{il}(y_l^+-x_i)
   -\sum_{l\in\mathcal J_i^-}a^-_{il}(y_l^--x_i).
\end{equation}
Following the stopped-target construction of drifting models
\citep{deng2026drifting}, the auxiliary objective is
\begin{equation}
  \mathcal L_{\mathrm{drift}}
  =\frac{1}{N_{\mathrm{act}}}\sum_i
  \left\|x_i-\operatorname{sg}\!\left(
    x_i+\eta\frac{\widehat V_i}{\|\widehat V_i\|_2+\epsilon}
  \right)\right\|_2^2,
\end{equation}
where $\eta$ controls the target displacement. The block-level variant applies
the same construction to concatenated normalized block features,
$b_r^S=k_r^{-1/2}\operatorname{vec}[x_{r,1};\ldots;x_{r,k_r}]$, and their
teacher counterparts, grouping blocks of the same length.

\subsection{Controlled comparisons}
The planned variants are summarized in Table~\ref{tab:variants}. We will keep
initialization, data, and training budgets comparable. During warm-up at fixed
$k=3$, we will sweep $\lambda_R\in[10^{-4},10^{-1}]$ and target an auxiliary-to-token
gradient-norm ratio of approximately $0.10$. We will then freeze the loss
coefficient and evaluate static block sizes $k\in\{2,3,4,5\}$, with an extended
sweep to $k=6,7$ if feasible. Training cost and memory will be reported alongside
inference performance.

\begin{table}[t]
  \caption{Planned objective variants and their diagnostic roles.}
  \label{tab:variants}
  \centering
  \begin{tabular}{lll}
    \toprule
    ID & Objective & Diagnostic role \\
    \midrule
    M0 & Hard on-policy MTP & Token-only baseline \\
    M1 & M0 + paired latent MSE & Pointwise alignment \\
    M2 & M0 + conditional full drift & Attraction and repulsion \\
    M3 & M0 + attraction-only drift & Repulsion ablation \\
    M4 & M0 + block-level full drift & Joint block coherence \\
    M5 & Soft on-policy KL distillation & Soft token supervision \\
    \bottomrule
  \end{tabular}
\end{table}

\subsection{Evaluation and decision rules}
Training will use mathematical reasoning traces from GSM8K and MetaMathQA.
The primary evaluations will cover GSM8K, AIME25, and GPQA Main; broader
capability checks will include BBH, IFEval, and CNN/DailyMail. We will report
task accuracy, instruction-following scores, and ROUGE-L as appropriate. The
frozen teacher under sequential decoding ($k=1$) provides the capability-retention
reference. All model-selection decisions will use held-out validation data.

We will measure token-level expected calibration error (ECE), following the
calibration perspective of~\citet{joshi2025calibration}. Separately, we will
compare teacher and student feature populations using sliced Wasserstein-2,
maximum mean discrepancy~\citep{gretton2012kernel}, Sinkhorn transport cost
\citep{cuturi2013sinkhorn}, and 1-nearest-neighbor two-sample accuracy, with
complementary distributional diagnostics motivated by
\citet{esteban2026kernel}. Features will use shared normalization and any PCA
projection fitted only on training features. Diagnostics will be reported per
horizon; nearest-neighbor accuracy approaching $50\%$ indicates reduced
population distinguishability under balanced sampling.

We will evaluate both static and confidence-adaptive decoding. Realized token
acceleration is
\begin{equation}
  k_{\mathrm{eff}}
  =\frac{\text{emitted non-padding tokens}}
        {\text{decoding forward passes}},
\end{equation}
which will be reported alongside measured throughput and end-to-end latency.
For a single operating-point summary, validation will select the fastest
setting within one percentage point of the sequential teacher's validation
accuracy, when such a setting exists; it will then be frozen for test evaluation.
We will report uncertainty over evaluation examples and, where feasible,
training seeds. If MSE outperforms drift, or token-only training outperforms
both regularizers, we will report that outcome and examine sampling noise,
regularization strength, and conditioning rather than assume the proposed
ordering holds.

\section{Risks}
\label{sec:risks}
\paragraph{Noisy or mismatched feature populations.}
Conditional drift depends on the size and relevance of local teacher and
student pools. Sparse horizon groups or unrelated contexts may produce noisy
updates. We will monitor pool sizes, kernel affinities, feature norms, and
per-horizon diagnostics; cap or subsample pools for tractability; and retain
paired MSE and attraction-only drift as diagnostic alternatives. Updates with
insufficient pool support will be skipped.

\paragraph{Overregularization and misleading alignment.}
Forcing feature alignment may reduce token accuracy or encourage feature
collapse. Matched auxiliary gradient strength and coefficient sweeps will help
separate the effect of the regularizer from loss scaling. Lower latent distances
or ECE alone will not be treated as evidence of improved reasoning: task scores
and the teacher-retention gap remain necessary checks.

\paragraph{Training overhead and deployment trade-offs.}
Pairwise kernel interactions increase training compute and memory, and a larger
configured block size need not yield proportional wall-clock speedup. We will
bound population sizes and profile teacher passes, kernel computation, and
training memory. Deployment comparisons will use identical hardware, batch
sizes, prompt lengths, and generation limits.

\paragraph{Generalization and evaluation validity.}
Training on mathematical traces may degrade other capabilities, motivating the
broader evaluation suite. Validation and test splits will remain separate, and
training data will be checked for benchmark overlap. Faster generation may also
amplify existing model errors; capability retention and confidence calibration
will be assessed before making practical deployment claims.

\section{Timeline \& Resources}
\label{sec:timeline-resources}
% Planning additions: the source manuscript specifies neither a calendar nor
% a hardware allocation. Durations and budgets below remain to be confirmed.
The work will proceed through four phases. Dates and durations remain to be
confirmed against the available compute allocation.
\begin{enumerate}
  \item \textbf{Baseline reproduction:} establish the on-policy MTP pipeline,
    sequential teacher reference, paired MSE baseline, and evaluation harness.
  \item \textbf{Drift implementation and pilot:} implement conditional pools,
    self-exclusion, and stopped-target updates; profile memory and runtime;
    calibrate regularization strength at $k=3$.
  \item \textbf{Main experiments and ablations:} run the block-size and
    confidence-threshold sweeps, compare all objective variants, and collect
    task, calibration, and latent-alignment measurements.
  \item \textbf{Analysis and reporting:} quantify uncertainty, analyze failure
    cases and capability retention, and prepare reproducible configurations
    and the final manuscript.
\end{enumerate}

Required resources include a pretrained autoregressive checkpoint, a frozen
teacher copy, GPUs with sufficient memory for student training and teacher
forward passes, checkpoint storage, and held-out evaluation data. A pilot run
will determine feasible batch and pool sizes and measure GPU-hours per run;
the total budget will then be estimated from the planned configurations and
seeds. The checkpoint, GPU type and count, storage budget, personnel allocation,
and phase durations are \textbf{TBD}.

\bibliographystyle{plainnat}
\bibliography{references}

\end{document}
